\begin{pastexam}[label={px:ntut-112-ee-ds-03}]{北科大電機所（戊組）112 學年度 第 3 題}
（追蹤）

Sort the following data in ascending order by using three methods:
\begin{center}700, 123, 487, 1101, 657, 53, 188, 25.\end{center}
Show your steps.\\
(a)Bubble Sort (5\%)\\
(b)Quick Sort (5\%)\\
(c)Radix Sort (5\%)\\
\hspace*{1em}[Note: Please choose the first key as pivot in quick sort.]
\end{pastexam}

\begin{pastsoln}
位置編號 $1$--$8$。

\textbf{(a) 氣泡排序}（資結書第~14.2~節、虛擬碼~14.1）：每一趟由左到右比較相鄰的兩個數，左邊比較大就交換，
一趟結束時這一趟範圍內最大的數到了最右邊（粗體是這一趟歸位的數）。
\begin{center}
\small
\begin{tabular}{lc*{8}{r}}
\toprule
趟 & 交換次數 & 1 & 2 & 3 & 4 & 5 & 6 & 7 & 8 \\
\midrule
原始 & & 700 & 123 & 487 & 1101 & 657 & 53 & 188 & 25 \\
1 & 6 & 123 & 487 & 700 & 657 & 53 & 188 & 25 & \textbf{1101} \\
2 & 4 & 123 & 487 & 657 & 53 & 188 & 25 & \textbf{700} & 1101 \\
3 & 3 & 123 & 487 & 53 & 188 & 25 & \textbf{657} & 700 & 1101 \\
4 & 3 & 123 & 53 & 188 & 25 & \textbf{487} & 657 & 700 & 1101 \\
5 & 2 & 53 & 123 & 25 & \textbf{188} & 487 & 657 & 700 & 1101 \\
6 & 1 & 53 & 25 & \textbf{123} & 188 & 487 & 657 & 700 & 1101 \\
7 & 1 & \textbf{25} & \textbf{53} & 123 & 188 & 487 & 657 & 700 & 1101 \\
\bottomrule
\end{tabular}
\end{center}
第 1 趟的細節：$700 > 123$ 換；$700 > 487$ 換；$700 < 1101$ 不換；$1101 > 657$、$1101 > 53$、$1101 > 188$、$1101 > 25$ 各換一次，
$1101$ 一路浮到最後。 最小的 $25$ 原本在最右邊，每趟只能往左移一格，所以要做滿 $n - 1 = 7$ 趟；交換總數
$6 + 4 + 3 + 3 + 2 + 1 + 1 = 20$，等於逆序對的個數。

\textbf{(b) 快速排序}（取第一個鍵當基準）。 題目只規定基準，沒有規定分割的細節，這裡用常見的「兩個指標從兩端
往中間找、交錯後把基準換到分界」的寫法（資結書第~14.5~節作答提醒提到的那一種）：$i$ 從 $p + 1$ 往右找
$\ge x$ 的數，$j$ 從 $r$ 往左找 $\le x$ 的數，$i < j$ 就交換、繼續；$i \ge j$ 時把基準 $A[p]$ 和 $A[j]$ 交換，基準
歸位在 $j$，遞迴 $[p..j-1]$ 和 $[j+1..r]$。

\begin{center}
\small
\begin{tabular}{cc*{8}{r}}
\toprule
範圍 & 基準 & 1 & 2 & 3 & 4 & 5 & 6 & 7 & 8 \\
\midrule
開始 & & 700 & 123 & 487 & 1101 & 657 & 53 & 188 & 25 \\
$[1..8]$ & 700 & 188 & 123 & 487 & 25 & 657 & 53 & \textbf{700} & 1101 \\
$[1..6]$ & 188 & 25 & 123 & 53 & \textbf{188} & 657 & 487 & 700 & 1101 \\
$[1..3]$ & 25 & \textbf{25} & 123 & 53 & 188 & 657 & 487 & 700 & 1101 \\
$[2..3]$ & 123 & 25 & 53 & \textbf{123} & 188 & 657 & 487 & 700 & 1101 \\
$[5..6]$ & 657 & 25 & 53 & 123 & 188 & 487 & \textbf{657} & 700 & 1101 \\
\bottomrule
\end{tabular}
\end{center}
每次分割的細節：
\begin{itemize}
\item $[1..8]$，$x = 700$：$i$ 跳過 $123, 487$ 停在 $1101$（位置 4），$j$ 停在 $25$（位置 8），交換得
$700, 123, 487, 25, 657, 53, 188, 1101$。 $i$ 再往右跳過 $657, 53, 188$ 停在 $1101$（位置 8），$j$ 停在 $188$（位置 7），
交錯，基準和 $A[7]$ 交換。
\item $[1..6] = 188, 123, 487, 25, 657, 53$，$x = 188$：$i$ 停在 $487$（位置 3），$j$ 停在 $53$（位置 6），交換得
$188, 123, 53, 25, 657, 487$；$i$ 停在 $657$（位置 5），$j$ 停在 $25$（位置 4），交錯，基準和 $A[4]$ 交換。
\item $[1..3] = 25, 123, 53$，$x = 25$：$i$ 停在 $123$，$j$ 一路往左停在基準自己（位置 1），基準不動。
\item $[2..3] = 123, 53$，$x = 123$：$i$ 走出範圍，$j$ 停在 $53$（位置 3），基準和 $A[3]$ 交換。
\item $[5..6] = 657, 487$：同理，交換成 $487, 657$。
\end{itemize}
只有一個數或空的範圍直接返回。 結果 $25, 53, 123, 188, 487, 657, 700, 1101$。

\textbf{(c) 基數排序}（LSD，資結書第~15.2.2~節）：以 $10$ 為基數，從個位數開始，每一趟依那一位分到 $0$--$9$
十個桶子，桶子之間依編號、桶子內依放入的先後收回（穩定）。 最大的 $1101$ 有 $4$ 位，做 $4$ 趟，位數不足的
前面補 $0$（例如 $25 = 0025$）。
\begin{center}
\small
\begin{tabular}{p{1.5cm}p{6.4cm}p{5.6cm}}
\toprule
趟（依哪一位） & 非空的桶子 & 收回後的順序 \\
\midrule
1（個位） & 0: 700；1: 1101；3: 123, 53；5: 25；7: 487, 657；8: 188 & 700, 1101, 123, 53, 25, 487, 657, 188 \\
2（十位） & 0: 700, 1101；2: 123, 25；5: 53, 657；8: 487, 188 & 700, 1101, 123, 25, 53, 657, 487, 188 \\
3（百位） & 0: 25, 53；1: 1101, 123, 188；4: 487；6: 657；7: 700 & 25, 53, 1101, 123, 188, 487, 657, 700 \\
4（千位） & 0: 25, 53, 123, 188, 487, 657, 700；1: 1101 & 25, 53, 123, 188, 487, 657, 700, 1101 \\
\bottomrule
\end{tabular}
\end{center}
每一趟都要照上一趟收回的順序放進桶子，例如第 3 趟百位是 $1$ 的依序是 $1101, 123, 188$（上一趟它們就是這個先後），
不能自己重新排。 第 4 趟後完成。
\end{pastsoln}
